\chapter{From arithmetic depth to a Gamma shape calculus}
\label{chap:shape-core}

Volumes IV--VI follow one continuous
route. A signed number-basis amplitude has a positive time-averaged Gram
kernel. Large depth fixes the shape and centering of a log-Gamma window. Exact
transforms of that window then produce a phase-space kernel, orthogonal
polynomial coordinates, shape refinement, and coherent smoothing. This chapter
introduces the shared objects once. Source references identify the original
volume and claim title. Elementary identities are rederived here. The
arithmetic asymptotic and classical orthogonal-polynomial inputs retain their
original hypotheses and evidence scope.

\section{The arithmetic seed and its positive lift}

Retain Chapter 1's signed transfer $C_{a,\beta,\chi}(k,t)$, with
$\beta>0$, $a>1$, and $\chi$ modulo $q$. Put $d_n=1-n^{-\beta}$.
The zero-depth convention retains $n=1$. We now add a two-depth positive kernel.
The positive companion is defined already for $a>1/2$:
\[
 \mathcal K_q(k,\ell;\tau)
 =\sum_{\substack{n\geq1\\(n,q)=1}}
 \mu(n)^2 n^{-2a-i\tau}d_n^{k+\ell}.
\]

\begin{proposition}[Time-average Gram lift]
For $a>1$, the long-time cross-correlation satisfies
\[
 \lim_{T\to\infty}\frac1{2T}\int_{-T}^{T}
 C_{a,\beta,\chi}(k,t+\tau)\overline{C_{a,\beta,\chi}(\ell,t)}\,dt
 =\mathcal K_q(k,\ell;\tau).
\]
For $a>1/2$, the kernel
$\mathcal K_q(k,\ell;t-s)$ is positive definite on depth-time pairs.
\end{proposition}

\begin{proof}
Absolute convergence for $a>1$ permits expansion of the correlation into a
double series and passage through the time average. A term with indices $n,m$
contains $e^{-it(\log n-\log m)}$. Its mean is one for $n=m$ and tends to zero
otherwise. This leaves the displayed squarefree diagonal.
For positivity, set
$v_{k,t}(n)=\mu(n)\chi(n)n^{-a}d_n^k n^{-it}$.
These vectors are square summable for $a>1/2$, and the kernel is
$\sum_n v_{k,t}(n)\overline{v_{\ell,s}(n)}$.
Every finite kernel matrix is a Gram matrix.
\end{proof}

This is Paper IV, ``Exact time-average identity'' and ``Positive-definite packet
kernel.'' Time averaging retains the positive diagonal power and removes the
off-diagonal frequency interference. It does not bound the original signed
amplitude at a particular time.

Define
\[
 \alpha=\frac{2a-1}{\beta},\qquad
 \delta_q=\frac1{\zeta(2)}\prod_{p\mid q}\frac{p}{p+1}.
\]
The fixed-modulus squarefree counting input is
$\sum_{n\leq X,(n,q)=1}\mu(n)^2=\delta_q X+O_q(X^{1/2})$.
Paper IV, ``Uniform complex packet asymptotic,'' uses partial summation to obtain
\[
 \mathcal K_q(k,\ell;\tau)
 \sim\frac{\delta_q}{\beta}\Gamma(\alpha+i\tau/\beta)
 (k+\ell)^{-\alpha-i\tau/\beta}.
\]
Paper V, ``Joint arithmetic phase-space limit,'' combines the diagonal
normalizations. Uniformly when $k/\ell$ stays in a compact subset of
$(0,\infty)$ and $\tau$ stays in a compact real interval, with
$h=\log(k/\ell)$ and $\eta=\tau/\beta$,
\begin{equation}
 e^{i\eta\log(2\sqrt{k\ell})}
 \frac{\mathcal K_q(k,\ell;\tau)}
 {\sqrt{\mathcal K_q(k,k;0)\mathcal K_q(\ell,\ell;0)}}
 \longrightarrow
 \frac{\Gamma(\alpha+i\eta)}{\Gamma(\alpha)}
 \operatorname{sech}^{\alpha+i\eta}(h/2).
 \label{eq:shape-arithmetic-limit}
\end{equation}
The positive domain is $a>1/2$. The signed series above and its absolute
time-average derivation use the stronger condition $a>1$.

\section{One window and its exact phase-space kernel}

For an independent shape $\alpha>0$, define
\[
 g_\alpha(y)=\Gamma(\alpha)^{-1/2}
 \exp\left(\frac{\alpha y}{2}-\frac{e^y}{2}\right),
 \qquad y\in\mathbb R.
\]
The substitution $u=e^y$ gives
$\int |g_\alpha|^2dy=\Gamma(\alpha)^{-1}
\int_0^\infty u^{\alpha-1}e^{-u}du=1$.
Thus $|g_\alpha|^2$ is a log-Gamma density. Its shape is fixed by the arithmetic
limit in \eqref{eq:shape-arithmetic-limit}, but the exact window theory makes
sense for every positive shape.

Use a second shape $\gamma$ to distinguish shape composition from the arithmetic
scale exponent. The centered cross-ambiguity convention is
\[
 A_{\alpha,\gamma}(h,\eta)
 =\int_{\mathbb R}g_\alpha(y+h/2)g_\gamma(y-h/2)e^{i\eta y}\,dy.
\]

\begin{theorem}[Cross-shape Gamma--sech identity]
For $\alpha,\gamma>0$ and real $h,\eta$, with $\lambda=(\alpha+\gamma)/2$,
\begin{equation}
 A_{\alpha,\gamma}(h,\eta)
 =e^{(\alpha-\gamma)h/4}
 \frac{\Gamma(\lambda+i\eta)}{\sqrt{\Gamma(\alpha)\Gamma(\gamma)}}
 \operatorname{sech}^{\lambda+i\eta}(h/2).
 \label{eq:shape-cross-ambiguity}
\end{equation}
\end{theorem}

\begin{proof}
Multiplication of the shifted windows leaves the constant prefactor
$e^{(\alpha-\gamma)h/4}/\sqrt{\Gamma(\alpha)\Gamma(\gamma)}$
and the integrand
$\exp((\lambda+i\eta)y-e^y\cosh(h/2))$.
The substitution $u=e^y\cosh(h/2)$ evaluates the remaining integral as
$\Gamma(\lambda+i\eta)\cosh^{-\lambda-i\eta}(h/2)$.
\end{proof}

This is Paper VI, ``Cross-shape Gamma--sech formula.'' At equal shapes it is
Paper V, ``Exact ambiguity formula,'' and precisely the arithmetic limit.
The Gamma function has no zeros and $\cosh(h/2)>0$ on the real line, so this
cross-ambiguity is nonzero on real phase space. At $h=\eta=0$ its value is
\[
 F(\alpha,\gamma)=
 \frac{\Gamma((\alpha+\gamma)/2)}{\sqrt{\Gamma(\alpha)\Gamma(\gamma)}}.
\]
Cauchy--Schwarz gives $0<F\leq1$, with equality only for equal shapes.
Indeed, the ratio of the two positive windows is a constant times
$e^{(\alpha-\gamma)y/2}$ and can be constant only when $\alpha=\gamma$.

\section{Fourier density and additive shape}

Under the unitary convention
$\widehat f(\omega)=(2\pi)^{-1/2}\int f(y)e^{-i\omega y}dy$, the same
Gamma integral gives
\[
 \widehat g_\alpha(\omega)=
 \frac{2^{\alpha/2-i\omega}\Gamma(\alpha/2-i\omega)}
 {\sqrt{2\pi\Gamma(\alpha)}}.
\]
Its squared modulus is the symmetric Meixner density
\[
 \rho_\alpha(\omega)=
 \frac{2^{\alpha-1}}{\pi\Gamma(\alpha)}
 |\Gamma(\alpha/2+i\omega)|^2.
\]
The Fourier autocorrelation identity and the zero-modulation ambiguity slice
give
\[
 \int_{\mathbb R}\rho_\alpha(\omega)e^{ih\omega}d\omega
 =\operatorname{sech}^{\alpha}(h/2).
\]
Consequently Fourier uniqueness proves the shape convolution law
\[
 \rho_{s+t}=\rho_s*\rho_t,\qquad s,t>0.
\]
These are Paper V, ``Fourier transform of the log-Gamma packet'' and ``Exact
characteristic function,'' followed by Paper VI, ``Meixner shape semigroup.''

The cross-ambiguity energy also separates exactly:
\begin{equation}
 \frac{|A_{\alpha,\gamma}(h,\eta)|^2}{2\pi}
 =p_{\alpha,\gamma}(h)\rho_{\alpha+\gamma}(\eta),
 \qquad
 p_{\alpha,\gamma}(h)=
 \frac{\Gamma(\alpha+\gamma)e^{(\alpha-\gamma)h/2}}
 {2^{\alpha+\gamma}\Gamma(\alpha)\Gamma(\gamma)}
 \operatorname{sech}^{\alpha+\gamma}(h/2).
 \label{eq:shape-factorization}
\end{equation}
Direct cancellation of constants proves the formula. Setting
$u=e^h/(1+e^h)$ transforms $p_{\alpha,\gamma}(h)dh$ into the
Beta$(\alpha,\gamma)$ probability law. Both factors are normalized. This is
Paper VI, ``Exact cross-spectrogram factorization.''

\section{Polynomial refinement and the common ladder}

The classical Meixner--Pollaczek polynomial input is specified by its generating
function, rather than by an implicit choice of sign or phase:
\[
 \sum_{n\geq0}P_{n,s}(x)z^n
 =(1-iz)^{-s/2+ix}(1+iz)^{-s/2-ix},\qquad |z|<1.
\]
Let $(s)_n=\Gamma(s+n)/\Gamma(s)$ and put
\[
 p_{n,s}(x)=\sqrt{\frac{n!}{(s)_n}}P_{n,s}(x).
\]
Classical orthogonality makes these a normalized polynomial basis in
$L^2(\rho_s(x)dx)$. The degree $n$ and shape $s$ are separate coordinates.

\begin{theorem}[Normalized two-child refinement]
For $\alpha,\gamma>0$,
\[
 p_{n,\alpha+\gamma}(x+y)
 =\sum_{r=0}^{n}
 \sqrt{\binom nr\frac{(\alpha)_r(\gamma)_{n-r}}
 {(\alpha+\gamma)_n}}
 p_{r,\alpha}(x)p_{n-r,\gamma}(y).
\]
The squared coefficients sum to one.
\end{theorem}

\begin{proof}
Multiplying the generating functions with shapes $\alpha,\gamma$ produces the
one with shape $\alpha+\gamma$ and argument $x+y$. Coefficient comparison gives
the unnormalized addition formula, and the displayed normalizations give the
square roots. The coefficient identity
$\sum_r\binom nr(\alpha)_r(\gamma)_{n-r}=(\alpha+\gamma)_n$
follows by comparing coefficients in
$(1-z)^{-\alpha}(1-z)^{-\gamma}=(1-z)^{-\alpha-\gamma}$.
\end{proof}

Paper VI, ``Two-child packet refinement,'' interprets the coefficient squares
as a beta-binomial law. For many children they become the
Dirichlet-multinomial coefficients
\[
 C_{n,\mathbf r}^{\mathbf s}
 =\sqrt{\frac{n!\prod_j(s_j)_{r_j}}
 {(\sum_js_j)_n\prod_jr_j!}},\qquad \sum_jr_j=n.
\]
Multiplication of generating functions proves this extension. Intermediate
Pochhammer factors cancel along every binary refinement tree, establishing the
coassociativity recorded in Paper VI.

There is a complementary position basis. Classical Laguerre orthogonality gives
\[
 L_{n,\alpha}^{\mathrm{packet}}(y)
 =\sqrt{\frac{n!}{\Gamma(n+\alpha)}}
 e^{\alpha y/2-e^y/2}L_n^{(\alpha-1)}(e^y).
\]
Paper V, ``Laguerre--Meixner Fourier duality,'' evaluates its transform as
\[
 \widehat{L_{n,\alpha}^{\mathrm{packet}}}(\omega)
 =i^n\widehat g_\alpha(\omega)p_{n,\alpha}(\omega).
\]
The phase $i^n$ is part of the identity. The polynomial frequency basis and the
Fourier-transformed position basis therefore differ by degree-dependent phases.

On an abstract orthonormal basis $e_n$, define
\[
 K_0e_n=(n+\alpha/2)e_n,\quad
 K_+e_n=\sqrt{(n+1)(n+\alpha)}e_{n+1},\quad
 K_-e_n=\sqrt{n(n+\alpha-1)}e_{n-1}.
\]
On its finite linear span, direct calculation gives
\[
 [K_0,K_+]=K_+,\quad [K_0,K_-]=-K_-,\quad
 [K_-,K_+]=2K_0,
\]
and the Casimir is $\frac\alpha2(\frac\alpha2-1)I$.
The Laguerre multiplication recurrence has matrix
$2K_0-K_+-K_-$; the real polynomial Meixner frequency recurrence has matrix
$(K_++K_-)/2$. These are two coordinate realizations of one ladder. A Fourier
comparison must retain the degree phases above. This is Paper V,
``$\mathfrak{su}(1,1)$ relations.''

\section{Shape transfer and its exact contraction spectrum}

Let $X_s$ have law $\rho_s$ and let $Y_t$ be an independent increment with law
$\rho_t$. Shape convolution gives $X_s+Y_t\sim\rho_{s+t}$. Conditioning the
refinement identity on $X_s$ removes all positive-degree increment polynomials
by orthogonality. Thus
\[
 \mathbb E[p_{n,s+t}(X_s+Y_t)\mid X_s]
 =\sigma_n(s,t)p_{n,s}(X_s),\qquad
 \sigma_n(s,t)=\sqrt{\frac{(s)_n}{(s+t)_n}}.
\]
For the transfer
$T_{s,t}:L^2(\rho_{s+t})\to L^2(\rho_s)$ given by
$T_{s,t}f(x)=\int f(x+y)\rho_t(y)dy$, these are its singular values in the
classical complete polynomial bases. They satisfy
\[
 \sigma_n(s,t)\sigma_n(s+t,u)=\sigma_n(s,t+u),\qquad
 \sigma_n(s,t)\sim
 \sqrt{\frac{\Gamma(s+t)}{\Gamma(s)}}n^{-t/2}.
\]
The transfer is compact for every $t>0$ and belongs to $S_p$ precisely when
$pt>2$. Its mean-zero operator norm is $\sqrt{s/(s+t)}$. These are Paper VI,
``Polynomial martingale law'' and ``Cocycle and compactness.'' The completeness
input matters for a statement about the entire operator, rather than only its
polynomial restriction.

\section{Coherent smoothing and its Gaussian limit}

Take the inner product conjugate-linear in its first argument. The centered
Weyl states are
\[
 \psi_{\alpha,u,\eta}(y)=e^{i\eta(y-u/2)}g_\alpha(y-u).
\]
Their exact overlap is
\[
 \langle\psi_{\alpha,u,\eta},\psi_{\gamma,v,\nu}\rangle
 =e^{i(\nu u-\eta v)/2}A_{\alpha,\gamma}(v-u,\nu-\eta).
\]
The Moyal identity for the normalized window gives the continuous tight frame
\[
 \frac1{2\pi}\int_{\mathbb R^2}
 |\langle\psi_{\alpha,u,\eta},f\rangle|^2du\,d\eta=\|f\|_2^2.
\]
Its coherent transform has a reproducing kernel equal to this overlap.
These are Paper VI, ``Cross-shape coherent overlap'' and ``Continuous tight
frame.''

The lower symbol of the associated anti-Wick operator is convolution with
\[
 B_\alpha(h,\eta)=\frac{|A_{\alpha,\alpha}(h,\eta)|^2}{2\pi}
 =p_\alpha(h)\rho_{2\alpha}(\eta),\qquad
 p_\alpha(h)=\frac{\Gamma(\alpha+1/2)}{2\sqrt\pi\Gamma(\alpha)}
 \operatorname{sech}^{2\alpha}(h/2).
\]
Its ordinary two-dimensional Fourier multiplier is
\[
 \widehat B_\alpha(\xi,x)
 =\frac{|\Gamma(\alpha+i\xi)|^2}{\Gamma(\alpha)^2}
 \operatorname{sech}^{2\alpha}(x/2).
\]
It is strictly positive on real dual space and tends to zero at infinity.
Therefore the Berezin convolution is injective on $L^2$ but has an unbounded
inverse on its range, as Paper VI states in ``Injectivity and ill-posed
inversion.'' Exact recoverability is not a stability bound.

For independent $H_\alpha\sim p_\alpha$ and
$\Omega_\alpha\sim\rho_{2\alpha}$, the standardized variables
\[
 Z_\alpha=\sqrt{\alpha/2}\,H_\alpha,\qquad
 W_\alpha=\sqrt{2/\alpha}\,\Omega_\alpha
\]
converge to independent standard Gaussians. For position this follows from the
Gamma-ratio characteristic function with argument
$\xi\sqrt{\alpha/2}$; for frequency it follows from
$\operatorname{sech}^{2\alpha}(x/\sqrt{2\alpha})\to e^{-x^2/2}$.
Paper VI, ``Gaussian marginal limits,'' proves this distributional limit.
Its stronger ``Heat-operator limit'' also uses controlled moments and
multiplier derivatives to assert Schwartz-seminorm convergence of the
rescaled convolution to $e^{\Delta/2}$. That stronger source claim is not
inferred here solely from convergence in distribution or a finite-grid plot.

\section{The signed branch and its retained boundary}

The positive shape calculus does not replace the signed arithmetic problem.
For $a>1$ the signed generating function is
\[
 \mathcal C_{a,\beta,\chi}(z,t)
 =\sum_{n\geq1}\frac{\mu(n)\chi(n)n^{-a-it}}{1-z(1-n^{-\beta})}.
\]
Paper IV continues it off the physical cut and proves
\[
 (1-z)\mathcal C_{a,\beta,\chi}(z,t)
 +z\mathcal C_{a+\beta,\beta,\chi}(z,t)=L(a+it,\chi)^{-1}.
\]
Its nonzero Hankel determinant supplies a unique rational Pad\'e matching
function with degree bounds $\deg P\leq m-1$, $\deg Q\leq m$ and
$Q(0)=1$, preserving the first $2m$ Taylor coefficients. An ordinary
$m$-term simple-packet decomposition additionally follows when $\deg Q=m$
and $Q$ has $m$ simple zeros. This chapter makes no universal claim of full
denominator degree, nonbreakdown, physical-interval signed nodes, or stable
inversion.

The principal signed coherence target and its exact RH equivalence are
stated once in Chapter 1. The shape windows, convolution laws, refinement
coefficients, and coherent frames provide positive structure; that open target
still concerns cancellation of the original M\"obius signs.
